% ===================================================================
%  اللوحة نمرة ٢ — جسم الإنسان. — الفسحة. الألعاب.
%
%  Traduction, faite en 2026, en arabe egyptien ecrit, sur l'ido de
%  texto/io. Regles de la colonne : voir 10-lawha-01.tex.
%
%  LE CORPS HUMAIN EST LE TABLEAU OU L'EGYPTIEN SE SEPARE LE PLUS DE
%  L'ARABE STANDARD, ET CE N'EST PAS UN CHOIX DE STYLE : ce sont les
%  mots que tout le monde emploie, et l'arabe standard en emploie
%  d'autres.
%
%    وش        le visage   MSA : وجه
%    مناخير    le nez      MSA : أنف
%    ودان      les oreilles MSA : أذنان
%    راس       la tete     MSA : رأس
%    دقن       le menton   MSA : ذقن
%    شنب       la moustache MSA : شارب
%    سوالف     les favoris MSA : سوالف الشعر
%    شفايف     les levres  MSA : شفاه
%    سنان      les dents   MSA : أسنان
%    رقبة      le cou      MSA : عنق
%    زور       la gorge    MSA : حلق
%    قفا       la nuque
%    ضهر       le dos      MSA : ظهر
%    كتاف      les epaules MSA : أكتاف
%    مصارين    les intestins MSA : أمعاء
%    دراع      le bras     MSA : ذراع
%    إيد       la main     MSA : يد
%    صوابع     les doigts  MSA : أصابع
%    ضوافر     les ongles  MSA : أظافر
%    كوع       le coude
%    بطة الرجل le mollet
%    عضم       les os      MSA : عظم
%
%  ON VOIT LA REGLE PHONOLOGIQUE A L'OEIL NU : le ذ passe a د (دقن,
%  دراع), le ظ passe a ض (ضهر, عضم), et le pluriel brise remplace le
%  pluriel savant (سنان, ضوافر, صوابع, كتاف, مصارين). Vingt-deux mots
%  d'un seul tableau, et pas un qui soit un emprunt : c'est de
%  l'arabe, mais ce n'est pas le meme.
%
%  LES JEUX, EUX, ONT LEURS NOMS PROPRES AU CAIRE :
%
%    العجلة (37)      la bicyclette   MSA : دراجة
%    الطيارة الورق (33) le cerf-volant
%    المرجيحة (35)    la balancoire
%    العوم (40)       la nage         MSA : السباحة
%    البلي (21)       les billes
%    النحلة (18)      la toupie
%    كورة (22, 36)    le ballon       MSA : كرة
%    نطّ (30)         sauter
%
%  ET DEUX JEUX QUE L'EGYPTIEN DISTINGUE EXACTEMENT COMME L'IDO :
%  « celo-trovo » (25) est الاستخبية, se cacher et se chercher ;
%  « blindo-ludo » (31) est الاستغماية, les yeux bandes. Le francais
%  de 1926 dit « cache-cache » et « colin-maillard » ; l'egyptien a
%  ses deux mots au meme endroit, sans qu'on ait eu a choisir.
% ===================================================================

%%K t02-apar-1 apar
\VUsaut{4.0mm}
\VUtitre{15.0pt}{60}{اللوحة نمرة ٢}
\VUsaut{2.0mm}
\VUtitre{11.4pt}{0}{جسم الإنسان. --- الفسحة.}
\VUtitre{11.4pt}{0}{الألعاب.}

%%K t02-tit-0 sub
\VUtitre{13.2pt}{0}{جسم الإنسان.}
\VUsaut{2.4mm}
\VUcentre{10.2pt}{0}{\textit{(درس بسيط في علوم الطبيعة.)}}
\VUsaut{3.0mm}

%%K t02-tit-1 sub pos=t02-01-1
\VUpk{10.2pt}{40}{الراس.}
\VUsaut{3.0mm}

%%K t02-01-1 p
1. --- أجزاء جسم الإنسان الرئيسية هي: \VUgras{الراس}، و\VUgras{الجدع}،
و\VUgras{الأطراف}.

%%K t02-02-1 p
2. --- الراس فيه جزئين: \VUgras{الجمجمة}، اللي جواها \VUgras{المخ} واللي
\VUgras{الشعر} مغطّيها، و\VUgras{الوش}.

%%K t02-03-1 p
3. --- في الرسمة بتاعتنا مش شايفين لا الجمجمة ولا المخ؛ بس شايفين بوضوح
أجزاء الوش المهمة.

%%K t02-04-1 p
4. --- الأول خالص \VUgras{الجبهة}، اللي بتدلّ على عقل قوي وتفكير شغّال على
طول. و\VUgras{الصدغين} واسعين أوي. و\VUgras{العين} حيّة ومفتوحة على الآخر.
و\VUgras{الحاجب} الكتيف و\VUgras{الرموش} الطويلة بيحموها.

%%K t02-05-1 p
5. --- وبعدين \VUgras{المناخير} و\VUgras{فتحات المناخير}. المناخير بتنفعنا نشمّ
ونتنفّس. وعلى كل ناحية من المناخير في \VUgras{الخدود}، وأبعد شوية
\VUgras{الودان}. وتحت الوش في \VUgras{الفم} و\VUgras{الدقن}.

%%K t02-06-1 p
6. --- إحنا مش شايفينهم كويس أوي، علشان \VUgras{الشنب} و\VUgras{السوالف}
مغطّيينهم علينا. بس إحنا عارفين إن الفم فيه \VUgras{الشفايف} الاتنين،
و\VUgras{الفك الأعلى} و\VUgras{الفك الأسفل} بـ\VUgras{سنانهم}، و\VUgras{اللسان}.
بالفم إحنا بنتكلّم وبناكل. وباللسان بندوق الأكل.

%%K t02-07-1 p
7. --- العين والودن والمناخير والفم، زيهم زي الصوابع، هم أعضاء الحواس
الخمسة: \VUgras{الشوف}، و\VUgras{السمع}، و\VUgras{الشمّ}، و\VUgras{الدوق}،
و\VUgras{اللمس}.

%%K t02-08-1 p
8. --- يبقى الراس واحد من أكتر أجزاء جسم الإنسان اللي تستاهل الاهتمام.

%%K t02-tit-2 sub
\VUsaut{4.4mm}
\VUpk{10.2pt}{40}{الجدع.}
\VUsaut{3.0mm}

%%K t02-09-1 p
9. --- \VUgras{الرقبة} بتوصّل الراس بالجدع. وفيها \VUgras{الزور}
و\VUgras{القفا}.

%%K t02-10-1 p
10. --- والجدع كمان فيه أعضاء أساسية كتير. في \VUgras{الصدر} في
\VUgras{الرئتين}، اللي بنتنفّس بيهم، و\VUgras{القلب} اللي هو مركز الحياة. ولما
القلب يبطّل ينبض، الكائن الحي بيموت.

%%K t02-11-1 p
11. --- و\VUgras{الضلوع} هي اللي شايلة الصدر.

%%K t02-12-1 p
12. --- وباقي أجزاء الجدع هي \VUgras{الجنب}، و\VUgras{الضهر}،
و\VUgras{الكتاف}، و\VUgras{البطن}. إحنا بننام على جنبنا أو على ضهرنا، وبنشيل
الحاجات التقيلة على كتافنا.

%%K t02-13-1 p
13. --- وجوّه البطن في \VUgras{المعدة} و\VUgras{المصارين}. دول بينفعونا نهضم
الأكل.

%%K t02-tit-3 sub
\VUsaut{4.4mm}
\VUpk{10.2pt}{40}{الأطراف.}
\VUsaut{3.0mm}

%%K t02-14-1 p
14. --- عندنا أربع \VUgras{أطراف}: \VUgras{دراعين} و\VUgras{رجلين}. بالدراعين
بناخد ونمسك ونرفع ونلمّ الحاجات؛ وبالرجلين بنقف ونمشي ونجري.

%%K t02-15-1 p
15. --- \VUgras{الكتف} بيوصّل الدراع بالجسم. و\VUgras{عضلة} قوية أوي، اللي هي
العضلة ذات الراسين، بتحرّك الدراع اللي \VUgras{الكوع} بيوصّله
بـ\VUgras{الساعد}. و\VUgras{الرسغ} هو مفصل \VUgras{الإيد}، اللي بتنتهي
بـ\VUgras{الصوابع} و\VUgras{الضوافر}. وعلى الإيد بنشوف \VUgras{عِرق} (وشريان)
اللي بيجري فيه \VUgras{الدم}.

%%K t02-16-1 p
16. --- وأجزاء الرجل هي: \VUgras{الفخد}، و\VUgras{الركبة} و\VUgras{باطن
الركبة}، و\VUgras{بطة الرجل} اللي \VUgras{لحمها} مغطّى بـ\VUgras{الجلد}،
و\VUgras{الكعب}، و\VUgras{القدم}. و\VUgras{عضم} الرجل تخين أوي وقوي أوي. وفي
آخر القدم في \VUgras{صوابع الرجل}. وباقي الأجزاء هي \VUgras{باطن القدم}
و\VUgras{الكعب الخلفي}.

%%K t02-17-1 p
17. --- وده تقريبًا الوصف الكامل لجسم الإنسان.

%%K t02-17-2 p
\textit{ملحوظة.} --- \textit{بمساعدة العبارة: «أنا بيوجعني... (راسي وهكذا)»
يقدر المدرِّس يستعمل الكلام ده كله تاني من ناحية} \VUgras{الحالة الصحية}
\textit{الكويسة أو الوحشة.}

%%K t02-tit-4 sub
\VUsaut{5.0mm}
\VUpk{10.2pt}{40}{الجمباز.}
\VUsaut{2.4mm}
\VUcentre{10.2pt}{0}{\textit{(بيحكي واحد من التلامذة.)}}
\VUsaut{3.0mm}

%%K t02-18-1 p
18. --- علشان نبقى ليّنين وخفاف وصحّتنا كويسة، لازم نمرّن رجلينا ودراعاتنا.
علشان كده بنلعب وبنعمل \VUgras{جمباز}.

%%K t02-19-1 p
19. --- في وسط الأسبوع التمرينين دول بيبقوا في حوش المدرسة (بُصّ اللوحة نمرة
١). إنما يوم الخميس ويوم الحد بنروح ملعب أخضر كبير، فيه مكان نجري فيه ونتسلّى.

%%K t02-20-1 p
20. --- أساتذتنا وأهالينا بيروحوا معانا هناك ساعات؛ إنما اللي بيدير التمارين
هو \VUgras{مدرِّس الجمباز}\textsuperscript{(12)}.

%%K t02-21-1 p
21. --- في الملعب، على شمال المدخل، في \VUgras{أجهزة الجمباز}\textsuperscript{(1)}.
إحنا بنتسلّق \VUgras{السلّم}\textsuperscript{(2)}، وبنتقلّب على
\VUgras{الحلق}\textsuperscript{(3)} وعلى \VUgras{التربيز}\textsuperscript{(4)}، وبنشدّ نفسنا بإيدينا لغاية فوق
\VUgras{سلّم الحبل}\textsuperscript{(5)} أو \VUgras{الحبل المعقّد}\textsuperscript{(6)}.

%%K t02-22-1 p
22. --- وفي نفس الوقت، زمايلنا بينطّوا فوق \VUgras{الحصان الخشب}\textsuperscript{(7)}
أو \VUgras{المتوازي}\textsuperscript{(11)}. وفي ناس تانية
\VUgras{بتتلاكم}\textsuperscript{(8)} أو \VUgras{بتتبارز}\textsuperscript{(9)} بالقناع وبـ\VUgras{الشيش}. وفي آخرين
\VUgras{بيرفعوا الأثقال}\textsuperscript{(10)}.

%%K t02-tit-5 sub
\VUsaut{4.6mm}
\VUpk{10.2pt}{40}{الألعاب.}
\VUsaut{3.0mm}

%%K t02-23-1 p
23. --- حوالين اللي بيعملوا جمباز، الولاد والبنات بيلعبوا \VUgras{ألعاب} مختلفة.
البنات بيتسلّوا بـ\VUgras{الطوق}\textsuperscript{(14)} بتاعهم؛ وبينطّوا
\VUgras{الحبل}\textsuperscript{(13)}، أو بيعزموا إخواتهم وولاد عمّهم على لعبة
\VUgras{الكروكيه}\textsuperscript{(15)}. وبيموتوا وهما بيطيّروا \VUgras{كورة} الخصم
بـ\VUgras{الشاكوش الخشب} بتاعهم، في اللحظة اللي هو فاكر إنه هيعدّي من القوس
بسهولة!

%%K t02-24-1 p
24. --- الولاد بيفضّلوا الألعاب اللي فيها عضلات أكتر. بيحبّوا يلعبوا
\VUgras{المخاريط}\textsuperscript{(16)}، وبيوقّعوها بمهارة بـ\VUgras{الكورة}\textsuperscript{(17)}.
وكمان ما بيرفضوش يلعبوا \VUgras{النحلة}\textsuperscript{(18)}
و\VUgras{البلبوكيه}\textsuperscript{(19)} أو حتى \VUgras{لعبة الضفدعة}
\textsuperscript{(20)}. إنما ألعابهم المفضّلة هي
\VUgras{البلي}\textsuperscript{(21)} و\VUgras{الكورة على الحيطة}
\textsuperscript{(22)}، علشان حيطة \VUgras{الصوبة}\textsuperscript{(26)} كويسة أوي علشان الكورة الخفيفة ترجع منها.

%%K t02-25-1 p
25. --- ويوحنا الصغير، وهو ماشي على \VUgras{العكاكيز}\textsuperscript{(24)} بتاعته،
شاطر أوي وعارف يحفظ توازنه كويس.

%%K t02-noto-1 noto
\VUnotes{143.2mm}{%
(*) ألعاب العيال والولاد ليها كتير أسامي وطنية بحتة. إحنا اجتهدنا نسمّيها
في الإيدو بأحسن ما نقدر، يا إما إننا نستوحي من الحركة نفسها اللي بتميّزها،
يا إما إننا نختار الاسم الوطني في البلد دي أو دي، لما ما يبقاش بعيد أوي عن
الصح. مثلًا: \VUgras{لعبة البرميل} (بالألماني والفرنساوي) هي \VUgras{لعبة
الضفدعة} \textsuperscript{(20)} (بالإنجليزي)، واللي فيها اللاعبين بيحاولوا يرموا
الأقراص الصغيرة المدوّرة من فم الضفدعة المعدنية اللي فاتحة بُقّها وواقفة على
الصندوق (البرميل) المخرّم.
\VUgras{النطّ على الكتاف}\textsuperscript{(30)} بيختلف عن \VUgras{النطّ على الضهر}
في حاجة واحدة بس: علشان ينطّوا فوق الراس، اللاعبين بيسندوا إيديهم على كتاف
زميلهم، بدل ما في «\VUgras{النطّ على الضهر}» بيسندوا إيديهم على ضهره بس. ---
أو كمان بعض المشتركين بيبقوا مايلين واللي باقيين هينطّوا فوق ضهرهم وهما
سانِدين إيديهم على الكتف، فالأولين لازم يستحمّلوا الخبطة من غير ما يقعوا،
والتانيين لازم ينطّوا من غير ما يفقدوا توازنهم (بُصّ اللوحة نفسها).%
}

%%K t02-25-1 p suite
وجنبه في كام زميل بيلعبوا \VUgras{الاستخبية}\textsuperscript{(25)} في الصوبة دي وورا
\VUgras{السياج}. وفي ناس تانية بيفضّلوا \VUgras{خمّن مين ضربك}\textsuperscript{(28)} أو
\VUgras{كرة الريشة}\textsuperscript{(29)} اللي بتطير من \VUgras{مضرب} للتاني.

%%K t02-26-1 p
26. --- في نصّ الملعب في شجرتين. وتحت واحدة منهم، سبعة أو تمن ولاد عملوا لعبة
\VUgras{النطّ على الكتاف}\textsuperscript{(30)} (*). اللعبة دي مش معروفة في كل
البلاد؛ إنما كل الناس عارفة إيه هو \VUgras{النطّ على الضهر}، اللي الولاد مش
بيلعبوه المرة دي.

%%K t02-tit-6 sub
\VUsaut{5.0mm}
\VUpk{10.2pt}{40}{الألعاب في الهوا الطلق.}
\VUsaut{3.4mm}

%%K t02-27-1 p
27. \VUgras{الاستغماية}\textsuperscript{(31)} كمان لعبة مسلّية أوي؛ البنات والولاد
بالدور بيحطّوا \VUgras{العصابة} على عينيهم ويمسكوا اللي مش شاطرين واللي بيسيبوا
نفسهم يتمسكوا.

%%K t02-28-1 p
28. --- في نصّ الملعب، أهي لعبة فرنساوية أوي بس فيها عنف شوية: دي
\VUgras{لعبة الدبّة}\textsuperscript{(32)}، وهي أكتر دوشة بكتير من لعبة

%%K t02-28-1 p suite
\VUgras{الطيارة الورق}\textsuperscript{(33)} الهادية أوي، وحتى من اللعبة الإنجليزية
\VUgras{التنس}\textsuperscript{(34)}.

%%K t02-29-1 p
29. --- الرياضة الأخيرة دي هي فرحة التلامذة الكبار. ومدرّسينا نفسهم بيلعبوها
بمزاجهم. وهي عايزة مهارة كبيرة مع خفّة، وأنا بحبّها أوي. وأنا بسيب للبنات
ألعاب \VUgras{المرجيحة}\textsuperscript{(35)}.

%%K t02-30-1 p
30. --- وأنا مش بحبّ لعبة إنجليزية تانية، اللي يمكن تبقى خطرة.

%%K t02-30-1b p
أنا قصدي \VUgras{كورة القدم}\textsuperscript{(36)}، اللي بتفرّح أخويا الكبير. أنا
بفضّل \VUgras{لعبة البارات}\textsuperscript{(38)} القديمة بتاعتنا، صحّية زيها بالظبط
وأقل وحشية بكتير.

%%K t02-tit-7 sub
\VUsaut{4.6mm}
\VUpk{10.2pt}{40}{العجلة والبالون.}
\VUsaut{3.0mm}

%%K t02-31-1 p
31. --- وكمان بحبّ ألفّ حوالين الملعب بـ\VUgras{العجلة}\textsuperscript{(37)}.

%%K t02-32-1 p
32. --- والسنة الجاية هتعلّم \VUgras{ركوب الخيل}\textsuperscript{(39)}. ما أحلى حصان
بيخطي بشياكة أو بيهرول، وبينطّ فوق الحواجز وهو بيجري!

%%K t02-33-1 p
33. --- وفي الصيف بنتعلّم \VUgras{العوم}\textsuperscript{(40)}. بنغطس ونستحمّى بمزاجنا
في ميّة النهر الصافية الساقعة. وساعات في الآخر بنطيّر \VUgras{بالون}\textsuperscript{(41)}
من ورق، بيطلع في الهوا بعظمة أول ما يتملّي بهوا سخن.

%%K t02-34-1 p
34. --- وفي كمان ألعاب تانية كتير، بس أنا مش هخلّص لو عدّيتها كلها؛ وأي حد
هيفهم من الملخّص ده إن يوم الخميس، الواحد مش بيزهق أوي في ملعبنا الحلو.

%%K t02-34-2 p
ملحوظة. --- \textit{المدرِّس هيقدر بسهولة يكمّل الفصل ده؛ ويوصف بتفصيل أكتر كل
لعبة من الألعاب المهمة.}
